\section{The remaining quadratic-eight rows at coefficient seven}
\label{sec:coefficient-seven-two-prime-procyclic}

\paragraph{The original coefficient and the two-prime distinction.}
Put
\[
 R=\mathbb Q(\zeta_7)^+,
 \qquad F=R\mathbb Q(\sqrt2),
 \qquad L=\mathbb Q(\zeta_7)\mathbb Q(\sqrt2).
 \tag{RP7.1}
\]
The real field $F$ has two primes above seven. We do not apply
the one-prime class-group lemma to this field. Instead we prove
directly that its maximal pro-$7$ extension unramified outside seven
has Galois group $\mathbb Z_7$, and that either local
decomposition group maps onto this quotient.

For a finite extension $E/\mathbb Q_7$, let $O$ be its original
ring of integers and put $D_O=E/O$. For a number field $M$ and
a discrete $7$-primary coefficient $A$, write $\mathcal S_7(M,A)$
for the full strict kernel: zero in $H^1(M_v,A)$ at every
$v\mid7$, zero in $H^1(I_v,A)$ at every other finite place,
and zero at real places. Real-place $H^1$ is zero on such a
coefficient. Finite localizations containing all coefficient
ramification primes give the same kernel.

Let $c_7$ be the ordinary cyclotomic character and let $\tau_7$
be its Teichm\"uller lift, with values in $\mu_6\subset\mathbb Z_7$.
Let $K=\mathbb Q(\omega)$ and $\eta_K=\chi_{-3}$, and use the
primitive quadratic character $\chi_8$ of $\mathbb Q(\sqrt2)$.

\begin{theorem}[AC-G201: original-integral descent through a two-prime field]
\label{AC-G201}
Let $\Gamma_F$ be the Galois group of the maximal extension of
$F$ unramified at finite primes outside seven. Its maximal
pro-$7$ quotient $P_F$ is canonically its cyclotomic quotient:
\[
 \begin{gathered}
 P_F\simeq\operatorname{Gal}(F\mathbb Q_\infty/F)
       \simeq\mathbb Z_7,\\
 G_{F_v}\longrightarrow P_F\ \hbox{is onto for each }v\mid7.
 \end{gathered}
 \tag{RP7.2}
\]
where $\mathbb Q_\infty/\mathbb Q$ is the cyclotomic
$\mathbb Z_7$-extension. Consequently, for every original $O$
and every continuous scalar $a:G_F\to1+\mathfrak m_O$
unramified outside seven,
\[
                     \mathcal S_7(F,D_O(a))=0.
 \tag{RP7.3}
\]

For every integer $e\geq1$ and $0\leq j<6$ with $j-e$ even,
put
\[
 a^{\rm nc}=\chi_8\tau_7^j c_7^{-e},
 \qquad a^{\rm cr}=a^{\rm nc}\eta_K.
 \tag{RP7.4}
\]
On this unchanged coefficient,
\[
 \begin{aligned}
 \mathcal S_7(\mathbb Q,D_O(a^{\rm nc}))&=0,\\
 \mathcal S_7(K,D_O(\chi_8\tau_7^j c_7^{-e}))
       &\simeq\mathcal S_7(\mathbb Q,D_O(a^{\rm cr})).
 \end{aligned}\tag{RP7.5}
\]
The second group is the critical rational component; it is not
asserted zero.

In the actual host and original coefficients of (FC.1)--(FC.2),
for $h\geq2$ and any integer $T\geq1$, the set
\[
 Y^{\rm two}_7(T)=
 \{(7,(\chi_8\tau_7^j)|_{G_K},e):
               1\leq e<T,\ 0\leq j<6,\ j-e\text{ even}\}
 \tag{RP7.6}
\]
consists of $3(T-1)$ distinct actual trivial-cubic triples.
All undisplayed finite factors are trivial. For every retained
set $J$ in (CP.3),
\[
                  \mathcal R_J=
                         \mathcal R_{J\setminus Y^{\rm two}_7(T)}.
 \tag{RP7.7}
\]
For $T\geq2$, the normalized original strict sum on any subset
of (RP7.6) is bounded by the single full critical aggregate
$\mathcal B_{\rm cr}(F_{\rm host},T)$ of AC-G191, with its same
original $b_7$ and polynomial cost. No positive-derived-index
splitting, common-window choice, or new error term is asserted.
\end{theorem}

\begin{proof}
\proofstep{1. Exact fields and completed class-quotient certificates.}
The real trace polynomial for $R$ is
$P_R(Y)=Y^3+Y^2-2Y-1$, with
$\Phi_7(Z)=Z^3P_R(Z+Z^{-1})$.
The fields $\mathbb Q(\zeta_7)$ and $\mathbb Q(\sqrt2)$ are
linearly disjoint: the former is unramified at two, whereas
the latter is ramified there. Thus $[L:\mathbb Q]=12$,
$[F:\mathbb Q]=6$, and $L/F$ is the quadratic CM extension.
Its involution $c$ fixes $\sqrt2$ and inverts $\zeta_7$.

The following exact resultants define these fields:
\[
 \begin{aligned}
 P_F(X)&=\operatorname{Res}_Y(P_R(Y),(X-Y)^2-2)\\
 &=X^6+2X^5-9X^4-14X^3+10X^2+8X+1,\\
 P_L(X)&=\operatorname{Res}_Y(\Phi_7(Y),(X-Y)^2-2)\\
 &=X^{12}+2X^{11}-9X^{10}-16X^9+43X^8+70X^7\\
 &\quad-97X^6-126X^5+211X^4+212X^3-37X^2-142X+127.
 \end{aligned}\tag{RP7.8}
\]
For either cyclotomic generator $\alpha$, the sum
$\alpha+\sqrt2$ is primitive. Indeed an automorphism fixing it
would have $\sigma(\alpha)-\alpha=2\sqrt2$ if it changed the
quadratic sign, contradicting disjointness; if it fixes that
sign, it must also fix $\alpha$. This proves the claimed
degrees and the field identifications in (RP7.8).

The saved fresh GP computation executes \texttt{bnfinit(P,1)}
and then \texttt{bnfcertify(B,1)} for each displayed polynomial.
The complete normal outputs give
\[
 \begin{array}{c|c|c|c|c}
 \text{field}&\text{signature}&\text{absolute discriminant}
                         &\text{computed quotient}&\text{flag one}\\\hline
 F&(6,0)&1229312&1&1\\
 L&(0,6)&74049191673856&C_2&1.
 \end{array}\tag{RP7.9}
\]
The applied PARI contract for flag one certifies that the
actual class group is a quotient of the computed group.
It does not certify a complete unit basis. In particular
\[
                    \operatorname{Cl}(L)[7]=0.
 \tag{RP7.10}
\]
The certificate for $F$ also gives its class group zero, but
the argument below only needs (RP7.10). No candidate class
number or unit-index assertion is assumed. Complete inputs,
outputs, successful exit receipt and binary data are retained.

\proofstep{2. Complex conjugation fixes both primes over seven.}
The two roots $\pm3$ of $X^2-2$ modulo seven are simple, so
Hensel lifting splits $\mathbb Q(\sqrt2)$ at seven.
The extension $\mathbb Q_7(\zeta_7)/\mathbb Q_7$ is totally
ramified of degree six: $\Phi_7(1+X)$ is Eisenstein at seven.
The fixed field of $\zeta_7\mapsto\zeta_7^{-1}$ has degree three.
Consequently $F$ has
two primes $v\mid7$, with
\[
 F_v\simeq\mathbb Q_7(\zeta_7)^+,
 \qquad L_w\simeq\mathbb Q_7(\zeta_7),
 \qquad [L_w:F_v]=2.
 \tag{RP7.11}
\]
There is one $w$ over each $v$. The CM involution $c$ therefore
fixes each of the two primes of $L$ over seven. This is the
local fact needed for the $S$-unit argument; it does not merge
the two distinct primes of $F$.

\proofstep{3. The minus $S$-unit quotient is at most one-dimensional.}
Let $S$ consist of the primes above seven and infinity, and set
$U=\mathcal O_{L,S}^{\times}$.
For every $u\in U$, the element $u/c(u)$ has valuation zero
at both primes over seven, since $c$ fixes each such prime.
It has valuation zero at every other finite prime by the
$S$-unit condition. It is therefore an ordinary algebraic unit.
For every complex embedding $\sigma$ of the CM field $L$,
$\sigma(c(u))=\overline{\sigma(u)}$; hence
\[
                  |\sigma(u/c(u))|=1.
 \tag{RP7.12}
\]
An algebraic integer all of whose conjugates have absolute
value one is a root of unity. For completeness, the monic
integer polynomials of all its positive powers have bounded
degree and bounded coefficients, by their elementary symmetric
functions. Only finitely many such polynomials, and therefore
finitely many roots, occur. Two powers coincide, proving the
assertion. We conclude that $(1-c)U\subset\mu(L)$.

The quotient $V=U/U^7$ is an $\mathbb F_7$-vector space, with
the natural action of $c$. Since two is invertible in
$\mathbb F_7$, every $v\in V^-$ satisfies
$v=\tfrac12(1-c)v$. Thus $V^-$ lies in the image of $\mu(L)$.
The finite root-of-unity group of a number field is cyclic,
so its quotient by seventh powers has dimension at most one.
Therefore
\[
                              \dim_{\mathbb F_7}V^-\leq1.
 \tag{RP7.13}
\]
No explicit unit basis, regulator determinant or generation
assertion for a proposed subgroup of units is needed.

\proofstep{4. Kummer theory with its exact conjugation twist.}
The extension $L/F$ is unramified at finite primes outside
seven, since it is obtained by adjoining $\zeta_7$.
Both maximal $S$-ramified extensions may therefore be taken
inside the same algebraic closure, giving the usual
index-two restriction sequence. Averaging over this index
two gives
\[
 H^1(G_{F,S},\mathbb F_7)
       \simeq H^1(G_{L,S},\mathbb F_7)^{c=1}.
 \tag{RP7.14}
\]
Choose a primitive seventh root to identify $\mathbb F_7$
with $\mu_7$ over $L$. Under conjugation this identification
changes sign, because $c$ acts on $\mu_7$ by inversion.
Thus the right side of (RP7.14) is
\[
                       H^1(G_{L,S},\mu_7)^{c=-1}.
 \tag{RP7.15}
\]
This sign includes the action on coefficients, not merely
conjugation on the Galois group.

The equivariant $S$-Kummer sequence is
\[
 0\longrightarrow U/U^7
 \longrightarrow H^1(G_{L,S},\mu_7)
 \longrightarrow\operatorname{Cl}(\mathcal O_{L,S})[7]
 \longrightarrow0.
 \tag{RP7.16}
\]
The last class group is a quotient of the ordinary class group.
It has no $7$-torsion by (RP7.10). Combining (RP7.13)--(RP7.16)
therefore proves
\[
                 \dim_{\mathbb F_7}H^1(G_{F,S},\mathbb F_7)
                    \leq1.
 \tag{RP7.17}
\]
The exact primary input is the complete $S$-Kummer sequence
and proof in Sharifi, Proposition~1.4.5; its naturality gives
the displayed equivariance. The CM unit argument has been
proved directly above, including its extra $S$-prime condition.

\proofstep{5. The pro-$7$ group and both local images.}
Constant-coefficient $H^1$ in (RP7.17) is
$\operatorname{Hom}_{\rm cont}(P_F,\mathbb F_7)$.
The real cyclotomic extension $\mathbb Q_\infty/\mathbb Q$
is unramified outside seven, and its degree-$7^n$ layers are
disjoint from $F$, whose degree is six. It gives a surjection
\[
                    P_F\longrightarrow\mathbb Z_7.
 \tag{RP7.18}
\]
In particular the dimension in (RP7.17) is exactly one.
These local cyclotomic layers are totally ramified: the
polynomial $\Phi_{7^{n+1}}(1+X)$ is Eisenstein, since modulo
seven it is a pure power of $X$ and its constant term is seven.

We give the elementary pro-group conclusion explicitly.
Every nontrivial finite quotient $Q$ of $P_F$ is a finite
$7$-group whose homomorphisms to $\mathbb F_7$ form a
one-dimensional space. Its maximal proper subgroups are
normal of index seven, so they are kernels of such homomorphisms.
Indeed, induct on the order through the nontrivial center:
a maximal subgroup either contains the center, allowing the
quotient induction, or multiplies with the center to the full
group, which makes it normal with a simple abelian quotient.
There is exactly one such maximal subgroup. Every element
outside it generates $Q$, since otherwise its generated
subgroup would lie in that maximal subgroup. Choose an
element of $P_F$ on which a nonzero continuous character is
nonzero. Its image generates every finite quotient, so its
closed cyclic subgroup is all of $P_F$.
An infinite procyclic pro-$7$ group is $\mathbb Z_7$.
The surjection (RP7.18) is then an isomorphism, proving the
first assertion of (RP7.2).

Locally, every $F_v/\mathbb Q_7$ in (RP7.11) has degree three.
It is disjoint from every degree-$7^n$ cyclotomic local layer.
Hence $G_{F_v}$ surjects onto the cyclotomic $\mathbb Z_7$
quotient. By the isomorphism just proved it surjects onto
$P_F$. This holds separately for each of the two primes.
It is an actual local-generation assertion, not a deduction
from generation by two unrelated local subgroups.

\proofstep{6. The full original local witness and descent.}
Let $a:G_F\to1+\mathfrak m_O$ satisfy the hypotheses of
(RP7.3), and let $z$ represent a strict cocycle with values
in $D_O(a)$. Its image is finite by compactness. Choose
one $v\mid7$. Full local zero supplies $x\in D_O$ with
$z(g)=(a(g)-1)x$ on $G_{F_v}$.
An original uniformizer power $\pi^m$ kills both this witness
and all values of $z$. The scalar action on $D_O[\pi^m]$
has finite $7$-group image $U_m$. The map
\[
 G_F\longrightarrow D_O[\pi^m]\rtimes U_m,
 \qquad g\longmapsto(z(g),a(g)|_{D_O[\pi^m]})
 \tag{RP7.19}
\]
is a homomorphism to a finite $7$-group. Outside seven inertia
acts trivially on the coefficient, so the zero inertia class
means $z$ itself vanishes there. The homomorphism factors
through $P_F$. By (RP7.2) every global image is represented
by a local element at the chosen $v$. The coboundary formula
on that local element uses the same witness $x$, and therefore
holds globally. This proves (RP7.3) on the original $O$.

Now suppose $j-e$ is even. The residual character of
$a^{\rm nc}$ in (RP7.4) is
$\chi_8\overline{\tau_7}^{j-e}$.
The quadratic factor is killed by $\mathbb Q(\sqrt2)$,
and the even power of $\tau_7$ factors through $R$.
Thus the restriction to $G_F$ is near one. It is fully
unramified outside seven: $\chi_8$ is already trivial
on $G_F$, and both $\tau_7$ and $c_7$ are unramified
at all other finite primes. A rational strict class
restricts to the zero group in (RP7.3). Restriction followed
by corestriction is multiplication by six, a unit on a
$7$-primary group. This proves the first line of (RP7.5)
without enlarging or replacing the coefficient ring.

The original induced lattice from $K/\mathbb Q$ splits into
the two characters $a^{\rm nc},a^{\rm cr}$ by
$(1\pm c)/2$, whose change-of-basis determinant is a unit
times two. Global and local Shapiro therefore identify the
$K$ strict group with the sum of the two rational strict
groups. At the split prime seven both local primes are retained;
at the ramified prime three the inertia quotient is $C_2$,
and inertia Shapiro uses its actual index-two subgroup.
At every other prime the unramified quadratic extension gives
the two or conjugate inertia summands. These are exactly
the full conditions in (FC.2). Extra unramified localization
primes do not change the kernel by inflation--restriction.
The first summand is zero, proving the second line of (RP7.5).

\proofstep{7. Exact actual labels and the unchanged critical charge.}
The host (FG.3) contains $\sqrt2$ and $\mathbb Q(\zeta_7)$;
all finite values of the characters in (RP7.6) lie already
in the specified original $O_7$. Their cubic and $q,d$
factors are trivial. For each weight exactly three of the
six values of $j$ have $j-e$ even. Distinct displayed labels
have distinct restrictions to $G_K$: their rational ratios
have conductor dividing $56$ and cannot be $\chi_{-3}$,
which is ramified at three. This proves $3(T-1)$, with
the empty interpretation at $T=1$.

For these labels $\chi_8\tau_7^j(-1)=(-1)^e$, so the
rational vanishing is precisely the noncritical summand
in (CP.3). This gives (RP7.7) term by term. The remaining
original strict groups in (RP7.5) are distinct members of
the same nonnegative critical aggregate in (CP.2), with
the same normalization $b_7$. Any subset is bounded by that
single aggregate; no factor three is introduced.
These rows are disjoint from the twenty-one-per-weight family
of AC-G197, exactly by the opposite parity condition on
$\chi_8$. A later window theorem must still use the same
window and its one existing total comparison/torsion cost.
\end{proof}

\paragraph{Attribution and source boundary.}
The parent proposed the two-prime field and the route through
a one-generator maximal pro-$7$ quotient, initially using
independent local unit images. The present author replaces
that regulator test by the exact CM $S$-unit/Kummer argument
in Steps 3--4 and supplies the complete proof and fresh
class-quotient computations. No black-box $p$-rationality
criterion, regulator generation, GRH or Iwasawa $\mu$-bound
is used. Foundational Kummer theory and the documented PARI
class-quotient certificate remain the explicitly stated inputs.
